\chapter{Elia}

\section{Overview}

This document summarises the key findings of a presentation delivered on 19 March 2026 by Rafael Feito Kiczak, Head of Scenarios, Market and Adequacy Studies at Elia, Belgium's Transmission System Operator (TSO). The presentation, entitled \textit{Challenges of the Energy Transition (seen through the eyes of a TSO)}, covers Elia's institutional role, the broader drivers of the energy transition, grid infrastructure challenges, flexibility considerations, and projected evolutions of the Belgian and European electricity mix through 2036 and beyond.

\section{Elia in a Nutshell}

Elia Group is the parent company of Belgium's national TSO and operates across both regulated and non-regulated activities.

\subsection{Regulated Entities}

\begin{itemize}[noitemsep]
    \item \textbf{Elia Belgium}: The national TSO managing on- and offshore transmission systems. It is 99.99\% owned by Elia Group and holds a monopolistic position in Belgium.
    \item \textbf{50Hertz}: The TSO for north-east Germany, 80\% owned by Elia Group, also holding a monopolistic position in its territory.
    \item \textbf{Nemo Link}: A 50/50 joint venture with National Grid (UK), providing a grid interconnection between Belgium and the United Kingdom.
\end{itemize}

\subsection{Non-Regulated Entities}

\begin{itemize}[noitemsep]
    \item \textbf{EGI (Elia Grid International)}: A 100\% subsidiary focused on international energy market consultancy and engineering services.
    \item \textbf{WindGrid}: A 100\% subsidiary focused on international offshore wind developments.
    \item \textbf{Energy Re}: A US-based entity in which Elia Group is a minority shareholder, acting as a vehicle to export technical expertise.
\end{itemize}

As a TSO, Elia fulfils four core societal tasks: managing the infrastructure, balancing the electricity system, facilitating the market, and acting as a trustee for the broader electricity system. It connects generation (renewable and conventional) to distribution networks, industrial clients, and consumers.

\section{The Energy Transition: Drivers and Scale}

\subsection{Drivers Beyond Climate}

The energy transition -- broadly understood as the shift away from fossil fuels -- is driven by a wide range of factors beyond climate ambition alone. These include public health considerations (e.g.\ reduction of local air pollution), regulatory obligations (e.g.\ Low Emission Zones), economic incentives (cheaper alternative technologies), energy security and independence, geopolitical motivations (reducing dependence on certain supplier countries), employment and welfare considerations, and the emergence of mature new technologies such as photovoltaic (PV) panels, electric vehicles (EVs), and heat pumps.

\subsection{The 25-Year Horizon}

The presentation frames the challenge with a striking symmetry: 25 years have passed since the euro was introduced in 2000, and 25 years remain until 2050, the target date for climate neutrality. During the first 25-year period, the EU achieved a reduction of approximately 37\% in CO\textsubscript{2} emissions and reached a 25\% share of renewable energy sources (RES). As of 2024, the EU electricity system features 47\% RES-E and 70\% low-carbon generation. The next 25 years must deliver far more dramatic changes, and they will not unfold in isolation from global geopolitical and technological disruptions.

\subsection{Key Systemic Trends}

Several concurrent trends are expected to further disrupt the energy system:

\begin{itemize}[noitemsep]
    \item \textbf{Falling technology costs}: Solar PV costs have dropped from \$0.74M/MW in 2015 to \$0.18M/MW in 2023; battery storage costs have similarly fallen sharply.
    \item \textbf{Nuclear renaissance}: A European nuclear alliance aims to achieve 150~GW of nuclear capacity by 2050.
    \item \textbf{Massive PV growth}: Global solar installations surpassed 60~GW per year for two consecutive years; the world is on track to install 29\% more solar capacity in 2024 than in 2023.
    \item \textbf{Critical material dependencies}: The transition requires large quantities of cobalt, copper, graphite, lithium, nickel, and rare earths, with significant supply chains concentrated in China.
    \item \textbf{Digitalisation}: Investment in data centres in the United States has surged exponentially, a trend mirrored globally.
    \item \textbf{Industrial price differentials}: European industrial electricity prices are roughly 158\% higher than those in the US or China, creating competitive pressure.
\end{itemize}

\subsection{Scale of Change by 2050}

Compared to today, the electricity system is expected by 2050 to see:

\begin{itemize}[noitemsep]
    \item Electricity consumption multiplied by \textbf{2}, driven by building heating, e-mobility, industrial electrification, and data centres.
    \item Installed renewables multiplied by \textbf{5 to 10}, spanning PV, onshore wind, and offshore wind.
    \item Installed storage multiplied by \textbf{10 to 20}, encompassing large-scale and residential batteries, and potentially thermal storage.
\end{itemize}

In Belgium specifically, electricity currently represents less than 20\% of total final energy consumption, and the country imports approximately 90\% of its energy needs in the form of fossil fuels. By 2050, electricity is expected to represent around 50\% of total energy needs. Electrification is identified as the most energy-efficient pathway to net zero: for example, using offshore wind electricity through a heat pump to provide space heating requires only 7~GW of installed capacity, compared to 36~GW if the same energy were routed through green hydrogen electrolysis.

\section{Elia's Network and Grid Congestion}

\subsection{The Belgian Grid}

As of 2022, Elia operates approximately 8,849~km of lines and cables in Belgium, spanning voltage levels from 30~kV to 380~kV, with 807 substations and converter stations. Grid planning is structured around a Federal Development Plan (FDP) covering the transmission grid (HVDC to 110~kV), updated every four years, and regional investment plans for the local transportation grid (30--70~kV), updated more frequently.

\subsection{Drivers of Grid Expansion}

Grid expansion is required in any future scenario. The horizontal (transmission) system must accommodate new loads and generation, requiring reinforcements onshore, offshore, and between countries. The vertical (distribution) system must adapt to handle decentralised generation and new loads, necessitating significant upgrades down to the street level.

Global investment in grid infrastructure, low-emission electricity, and end-use technologies is set to roughly double this decade compared to pre-2020 levels, according to IEA data. Belgium alone plans \euro8.7~billion in grid investments between 2024 and 2028.

\subsection{The Congestion Problem}

The volume of new connection requests is surging far beyond any prior forecast or approved plans:

\begin{itemize}[noitemsep]
    \item \textbf{Large-scale batteries}: Existing capacity is 0.3~GW; requests in process total 30~GW (\textasciitilde10x the pipeline from approved plans).
    \item \textbf{Data centres}: Existing demand is 3~TWh; requests in process total 70~TWh (\textasciitilde8x the pipeline).
    \item \textbf{Electrification}: Current commitments point to 123~TWh of additional load; reserved and in-process requests reach 153~TWh (+33\%).
\end{itemize}

In total, approximately 56~GW of capacity is currently in the connection study queue, representing roughly 4.5 times Belgium's peak electricity demand. This phenomenon is observed globally: the UK has 716~GW in its connection queue (15x peak load), Italy has 702~GW (12x), and similar dynamics are visible in Germany, the Netherlands, and US system operators such as MISO and ERCOT.

New transmission grid connection studies requested have grown from 45 in 2020 to a projected \textasciitilde280 in 2025 (a more than sixfold increase). DSO-connected requests have grown from 359 to \textasciitilde1,700 over the same period.

\subsection{Challenges in Grid Development}

Key structural challenges include: uncertainty over the future generation and consumption mix (compounded by very long infrastructure deployment times of 9--12 years for high-voltage lines); surging volumes of grid connection studies and capacity reservations; increased stakeholder interest and scrutiny; the need for closer TSO--DSO collaboration; and a fragmented regulatory framework involving multiple federal and regional authorities.

\subsection{Responses to Congestion}

Elia identifies five priority action areas to address grid congestion:

\begin{enumerate}[noitemsep]
    \item \textbf{More study capacity}: Scaling connection study throughput to handle the surge in requests and issuing faster, more transparent study decisions.
    \item \textbf{Accelerate CAPEX delivery}: Converting ordered connections into physical assets more quickly.
    \item \textbf{Societal choices}: Deciding for which types and volumes of capacity the grid should be prepared, and steering this in line with industrial and societal goals.
    \item \textbf{Active queue management}: Moving away from a first-come, first-served approach to prevent speculative or duplicative projects from crowding out real projects.
    \item \textbf{Unlocking flexibility at scale}: Using demand flexibility to reduce peak stress on the grid and limit the total need for physical reinforcements, both at high-voltage and distribution levels.
\end{enumerate}

\section{Flexibility}

\subsection{Why Flexibility Matters}

As renewable energy penetration increases, so does price volatility. The day-ahead price spread in Belgium in June has risen from approximately \euro30/MW in 2021 to \euro150/MW in 2024. On certain days, imbalance prices swing between +\euro350/MWh and -\euro650/MWh within a single day. Over 400 hours of negative prices were recorded in Belgium in 2024 alone. Europe surpassed 60~GW in annual PV installations for two consecutive years, driving record negative price occurrences across Western Europe.

By 2036, RES is expected to set the marginal electricity price in over 20\% of hours, with ``other technologies'' (hydro, storage, demand flexibility) setting the price in 46\% of hours, displacing gas as the dominant price-setter.

\subsection{Flexibility Across Time Frames}

Flexibility requirements span multiple time horizons:

\begin{itemize}[noitemsep]
    \item \textbf{Long term}: Captured in adequacy and economic studies, addressing structural capacity needs.
    \item \textbf{Day-ahead to real-time}: Covered in flexibility studies, addressing portfolio balancing, intra-day forecast updates, and residual imbalances.
\end{itemize}

\subsection{Electric Vehicles as a Flexibility Resource}

A typical EV has a 50~kWh battery but consumes only around 10~kWh per day on average, implying a large reservoir of flexibility. For 1 million EVs, optimal spread charging would require only 300~MW of additional grid capacity; today's reality produces 600--700~MW; and worst-case uncoordinated charging could require 7,000~MW. This illustrates both the risk of unmanaged EV uptake and the opportunity of smart charging.

E-mobility is expected to represent around 15~TWh of electricity demand in Belgium by 2035 (up from 1~TWh today), driven primarily by company cars in the near term and subsequently by trucks and buses.

Local regulatory incentives play a significant role in shaping charging profiles. In Flanders, capacity tariffs encourage flatter load profiles; in Wallonia, time-of-use tariffs create different patterns. As of September 2025, approximately 20,552 dynamic contracts (linked to real-time market prices) exist in Belgium, with 99.5\% concentrated in Flanders. The number of such contracts is expected to rise sharply. Unlocking flexibility at scale requires simultaneously working on asset volume, asset flexibility capability, control signals, control capability, appropriate metering, and customer engagement.

\section{Energy Mix: Past and Future}

\subsection{Belgian Electricity Mix}

Belgium's electricity generation has historically been dominated by nuclear power. Coal has been phased out, and gas-fired generation has partially filled the gap. As of 2024, renewables represent 31\% of the Belgian electricity mix. Looking forward, the mix will increasingly include offshore wind, onshore wind, and solar PV, while nuclear's contribution depends on policy choices. Gas is projected to remain a significant balancing source in the near and medium term, covering 29--37\% of demand by 2030 before declining to 18--22\% by 2036. Net imports are projected to rise from 12\% of demand today to 19--22\% by 2036.

At the European level, RES is expected to represent over 70\% of electricity generation by 2035, with solar and offshore wind as the main growth drivers.

\subsection{Belgium's Energy Gap}

Even with all currently approved policies in place -- including 4.4~GW of offshore wind, the 10-year extension of Doel 4 and Tihange 3, and regional renewable plans -- total domestic electricity supply will fall substantially short of projected demand. The Belgian energy gap in 2036 is estimated at 60--70~TWh, narrowing to 45--55~TWh if 2~GW of new nuclear is added, and to 38--48~TWh with 3~GW of new nuclear. By 2050, the residual gap on top of the central domestic RES supply scenario is estimated at 70--90~TWh, which must be filled through some combination of additional renewables, new nuclear, imports, or sufficiency measures.

\subsection{Five Strategic Challenges}

The presentation closes by identifying five overarching strategic challenges for the Belgian electricity system:

\begin{enumerate}[noitemsep]
    \item \textbf{Electrification}: Progressing but at uneven pace across segments -- delays in large industrial projects, but higher-than-expected loads at the distribution level from e-mobility and small industries.
    \item \textbf{RES deployment}: Residential PV growing faster than anticipated, but uncertainty and delays in offshore development.
    \item \textbf{Data centres}: Connection requests and projections driven by AI demand far exceed previous planning assumptions; 13~TWh currently reserved, 55~TWh in study.
    \item \textbf{Large-scale batteries}: Similarly, 9~GW reserved, 21~GW in study -- well beyond earlier plans.
    \item \textbf{Nuclear policy shift}: Government signals now point toward new nuclear capacity and further plant life extensions, while current grid and adequacy plans only account for the Doel 4/Tihange 3 extension.
\end{enumerate}

These five challenges collectively require societal choices about reasonable ambitions, clear policy signals, and coordinated action across TSOs, DSOs, regulators, governments, and market participants. The energy transition is not only a technical challenge -- it is a governance, regulatory, and societal one.

\section{Podcast -- Raphaël Fetokitschak, Elia}

\subsubsection{Quelle est la fonction d'Elia ? Est-ce que Elia ne gère que des centrales ?}

Elia est le gestionnaire du réseau de transport d'électricité à haute tension en Belgique. Son rôle principal est de faire circuler les électrons depuis les points de production vers les réseaux de distribution, en opérant et en construisant les grandes lignes à haute tension visibles sur les autoroutes. Elia ne gère pas uniquement des centrales classiques : elle supervise également la production décentralisée, notamment les panneaux photovoltaïques (plus de 11~GW installés) et l'éolien offshore (2,3~GW) et onshore (plus de 2~GW). Au-delà du transport physique, Elia assure l'équilibre en temps réel entre production et consommation, en maintenant la fréquence du réseau à 50~Hz, et facilite certains marchés, notamment les services auxiliaires nécessaires à l'opération du réseau.

\subsubsection{Comment est-ce que Elia va évoluer ?}

Elia doit évoluer pour accompagner l'électrification croissante de l'économie belge. Aujourd'hui, l'électricité ne représente qu'environ 20~\% de la consommation finale d'énergie en Belgique (environ 80~TWh), mais les projections anticipent un doublement, voire un triplement de cette consommation, sous l'effet de l'électrification des transports, du chauffage résidentiel (via les pompes à chaleur) et des usages industriels. Cela implique des investissements massifs : environ 7,5~milliards d'euros sont planifiés pour les quatre prochaines années. Les délais de réalisation des projets de réseau étant de 10 à 15~ans, Elia se projette dès aujourd'hui à l'horizon 2040--2050, voire 2060, pour anticiper les besoins futurs en infrastructures.

\subsubsection{Comment est-ce que l'on se connecte au réseau de transport ? Quelle est la procédure ? Qui peut y être connecté ? Pourquoi est-ce que la demande de connexions a augmenté ?}

Les acteurs dont la puissance dépasse 25~MW peuvent se raccorder au réseau de transport géré par Elia, qui opère des tensions allant de 30~kV à 380~kV. La procédure passe d'abord par une étude d'orientation, puis par une étude de détail, débouchant sur une réservation de capacité et enfin sur la signature d'un contrat de raccordement. Le nombre de demandes a explosé ces dernières années~-- phénomène mondial et non belgo-belge~-- en raison de l'afflux de projets de batteries, de data centers et d'autres usages liés à l'électrification. Elia est passée d'une trentaine de dossiers traités annuellement à environ 150 études de raccordement par an. Ce pic s'explique aussi par un système de file d'attente en mode << premier arrivé, premier servi >>, avec des conditions d'entrée peu contraignantes et des coûts de dépôt faibles, favorisant la spéculation.

\subsubsection{Quelle est la capacité de batterie réservée pour les batteries ? Qu'est-ce que c'est vis-à-vis de la pointe de la demande ?}

Environ 10~GW de capacité sont aujourd'hui réservés sur le réseau d'Elia pour des projets de batteries. À titre de comparaison, la pointe de consommation hivernale en Belgique se situe entre 12 et 13~GW, et la charge observée en temps réel tourne autour de 7~GW en dehors des pics. Cette capacité réservée pour les batteries est donc du même ordre de grandeur que la pointe nationale, ce qui illustre l'ampleur des demandes reçues, bien au-delà des hypothèses initiales des plans de développement qui prévoyaient environ 3~GW pour 2030--2035.

\subsubsection{Pourquoi est-ce qu'il y a une spéculation sur ces capacités ?}

La spéculation s'explique principalement par la faiblesse des barrières à l'entrée du processus de raccordement : il n'est pas nécessaire de disposer d'un terrain, ni d'un projet finalisé pour déposer une demande. Les coûts associés à cette démarche sont très faibles, ce qui incite de nombreux acteurs à réserver de la capacité sans certitude quant à la réalisation effective du projet. Le mécanisme du premier arrivé, premier servi a amplifié cet engouement, conduisant à une file d'attente disproportionnée. Ce phénomène est observé à l'échelle mondiale, notamment en Allemagne (plus de 200~GW de batteries en demande de raccordement) et au Royaume-Uni.

\subsubsection{Concrètement, qu'est-ce que l'on doit fournir à la personne qui demande à se raccorder ?}

L'étude de raccordement fournie au demandeur précise à quel poste du réseau le raccordement est envisageable, selon quelles modalités techniques, et propose une estimation du coût associé. Le niveau de détail varie selon le stade de la procédure~: une étude d'orientation donne une première indication, tandis qu'une étude de détail affine les paramètres techniques et financiers. Ces études sont réalisées par une équipe dédiée au sein du département de développement du réseau, dont les effectifs ont triplé pour faire face à l'augmentation des demandes.

\subsubsection{Dans la préparation de l'épisode, on a parlé du paradoxe de l'électrification. Quel est-il ?}

L'électrification est reconnue comme la voie la plus économique et la plus efficace pour atteindre les objectifs climatiques. Cependant, elle nécessite d'importants investissements en capital (\textit{capex}) dans les infrastructures de réseau. Ces investissements se répercutent sur les tarifs de transport et de distribution, ce qui risque de renchérir le coût final de l'électricité. Or, un prix de l'électricité plus élevé peut freiner précisément l'électrification que l'on cherche à accélérer~-- notamment pour les ménages hésitant entre une pompe à chaleur et une chaudière au gaz, ou pour les industries exposées à la concurrence internationale. Ce cercle potentiellement vicieux constitue le paradoxe de l'électrification, qui appelle des mesures politiques correctrices, telles que des tax shifts entre le gaz et l'électricité ou des soutiens ciblés aux consommateurs industriels énergivores.

\subsubsection{Avant d'arriver à l'adéquation, pourriez-vous nous parler du scénario de blackout ? Comment on établit les normes pour rendre un réseau adéquat ?}

Un blackout correspond à un arrêt complet du système électrique, entraînant la paralysie de l'ensemble des services qui en dépendent~: eau, télécommunications, distribution de carburant, commerces. La panne espagnole de 2025 en est un exemple récent. L'adéquation, telle qu'Elia la modélise, ne cherche pas à prévenir les blackouts opérationnels~-- qui relèvent d'une autre expertise~-- mais à s'assurer que, sur un horizon de 1 à 10~ans, le système dispose d'une capacité de production, de flexibilité et de stockage suffisante pour répondre à la demande. La Belgique est modélisée comme une \og plaque de cuivre \fg{}, c'est-à-dire sans tenir compte des contraintes internes du réseau, en se concentrant sur l'équilibre offre-demande heure par heure. La norme de fiabilité retenue est le \textit{Loss of Load Expectation} (LOLE) de 3~heures en moyenne, défini par la loi belge. Cette norme est calculée de manière probabiliste sur 600 années simulées pour une année cible donnée, en faisant varier les conditions climatiques, les arrêts fortuits de centrales et les capacités d'interconnexion européennes. Elle est calibrée en rapportant la valeur de l'énergie non fournie (\textit{value of lost load}) au coût d'un nouvel entrant sur le marché de la production.

\subsubsection{Donc on ne s'intéresse pas à la structure du réseau, mais plutôt à la production et consommation du réseau ? C'est correct ?}

C'est en effet la logique des études d'adéquation : on s'assure qu'à chaque heure, l'offre de production disponible est suffisante pour couvrir la demande, indépendamment de la topologie du réseau. La modélisation en plaque de cuivre suppose que l'électricité peut circuler librement entre tous les points du système. La question de savoir si la bonne ligne est au bon endroit pour acheminer physiquement cette énergie relève d'analyses plus proches du temps réel, qui ne sont pas l'objet des études d'adéquation à long terme.

\subsubsection{Sur les résultats probabilistes, est-ce la moyenne qui permet de juger que le réseau est acceptable, ou utilise-t-on un autre critère probabiliste, par exemple le P95 ?}

La norme belge actuelle repose sur la moyenne des 600 années simulées, fixée à 3~heures de LOLE. Il est possible qu'une année simulée affiche 50~heures de défaillance et une autre zéro~: seule la moyenne est réglementairement contraignante. Entre 2013 et 2018 ou 2019, la Belgique appliquait un double critère~: les 3~heures en moyenne et un critère complémentaire en P95 (probabilité de ne pas dépasser 20~heures une année sur vingt). Ce double critère a été abandonné car il n'était pas aligné avec la réglementation européenne issue du \textit{Clean Energy Package}, qui a harmonisé les méthodologies de calcul de la norme de fiabilité entre États membres. Le simulateur utilisé, développé conjointement avec RTE, est le modèle open source Antares.

\subsubsection{Ceci est un brown-out dans ce cas, c'est correct ?}

Oui. Les situations de pénurie modélisées dans les études d'adéquation correspondent à des délestages partiels et contrôlés~-- volontaires ou imposés~-- d'une partie de la charge du réseau, et non à un arrêt complet du système. On parle effectivement de brown-out plutôt que de blackout.

\subsubsection{Il y a aussi la quantité d'énergie non servie, c'est correct ?}

Oui. En complément du LOLE, qui mesure le nombre d'heures de défaillance, les études quantifient également l'\textit{Energy Not Served} (ENS), c'est-à-dire la quantité d'énergie qui n'a pas pu être fournie durant ces heures de pénurie. Cet indicateur permet d'évaluer la gravité des épisodes de défaillance au-delà de leur simple fréquence.

\subsubsection{Est-ce que Entsoe centralise ces critères ?}

ENTSO-E produit chaque année une étude européenne d'adéquation, l'\textit{European Resource Adequacy Assessment} (ERAA), à laquelle Elia contribue avec la casquette de gestionnaire de réseau belge. Cette démarche a succédé aux anciens \textit{Systems Outlook and Adequacy Forecasts}, réalisés dans des tableurs et sans véritable modélisation probabiliste des interconnexions. Une première étude coordonnée avait été lancée vers 2014--2015 dans le cadre du Forum Pentalatéral de l'Énergie (Belgique, Luxembourg, France, Allemagne, Pays-Bas, Suisse, Autriche). L'ERAA est désormais inscrit dans la réglementation européenne, avec une collecte harmonisée des hypothèses auprès de chaque gestionnaire de réseau. Chaque pays souhaitant mettre en place un mécanisme de capacité doit au préalable définir sa norme de fiabilité selon une méthodologie réglementée et la notifier à la Commission européenne.

\subsubsection{Les rapports sont mis à jour pour la cohérence entre les acteurs ? Ou pour prendre en compte les changements d'hypothèse ?}

Les deux raisons sont valables. Les hypothèses évoluent très rapidement~-- prolongation ou non du nucléaire, fermeture du charbon en Allemagne, nouvelles politiques industrielles~-- de sorte qu'entre la photo prise en début d'année et la publication du rapport en fin d'année, des changements majeurs peuvent survenir. La mise à jour vise donc à la fois à maintenir la cohérence entre les gestionnaires de réseau européens, qui échangent en bilatéral et au sein d'ENTSO-E, et à intégrer les révisions d'hypothèses les plus récentes. En Belgique, Elia réalise son étude d'adéquation et de flexibilité tous les deux ans, avec une consultation publique en novembre ouverte aux acteurs de marché, aux régulateurs et à la Commission européenne.

\subsubsection{Pourriez-vous nous parler des précédents rapports que vous avez produits ?}

La première étude d'adéquation belge a été réalisée en 2016, à la demande de la ministre Marghem. Depuis, cinq éditions ont été publiées. Le contexte ayant motivé ces travaux remonte à la loi de 2003 sur la sortie du nucléaire, qui prévoyait la fermeture progressive des réacteurs à partir de 2015, combinée à l'arrêt de plusieurs centrales au charbon (Royen, Langerlo) et aux difficultés économiques des centrales au gaz après la crise de 2008. Face au risque d'inadéquation, une réserve stratégique avait été mise en place pour maintenir hors marché des centrales économiquement non viables mais utiles à la sécurité d'approvisionnement.

\subsubsection{Auriez-vous des anecdotes sur ces travaux ?}

Deux épisodes marquants sont évoqués. En 2014, le sabotage de la turbine à vapeur de Doel~4, conjugué aux arrêts de Doel~3 et Tihange~2 pour des inclusions d'hydrogène, a contraint l'équipe à recalculer en urgence les risques pour l'hiver 2014--2015, en grande partie pendant les vacances d'été. En 2018, la découverte de défauts dans le béton des bâtiments annexes de plusieurs réacteurs belges a provoqué leur mise à l'arrêt successive, plongeant le système dans une situation de pré-crise qui a nécessité la mise en place de groupes de travail d'urgence avec le régulateur et l'administration. Ces événements illustrent la nature générique des défauts affectant les parcs nucléaires construits en série~-- comme la corrosion sous contrainte identifiée en France en 2021--2022~--, et la nécessité de modéliser ces risques dans les études d'adéquation via des scénarios de sensibilité.

\subsubsection{Votre rapport de 2021 prévoyait 30~TWh d'import, est-ce adéquat de faire de si grands imports ? Comment en vient-on à ce résultat ? Quelle fut la réaction face à ce résultat ?}

Les 30~TWh d'import prévus en 2021 sont le résultat mécanique de la fermeture du parc nucléaire belge, qui produisait alors plus de 40~TWh, soit environ la moitié du mix électrique. Ces TWh devant être remplacés, une partie l'est par du gaz, l'autre par des importations. Ce résultat est purement économique~: si de l'électricité moins chère est disponible à l'étranger (nucléaire français, éolien allemand), le marché l'importe naturellement. Cela ne signifie pas que le système n'est pas adéquat~: l'adéquation porte sur la capacité des pays voisins à fournir de l'électricité à la Belgique lors des moments critiques (froid hivernal, absence de vent), et non sur le solde énergétique annuel. Par ailleurs, Elia ne perçoit aucun avantage financier lié au volume importé ou exporté. La Belgique importe déjà 80 à 90~\% de ses besoins énergétiques totaux (pétrole, gaz), les importations d'électricité s'inscrivent donc dans cette réalité structurelle.

\subsubsection{Elia ne gagne pas plus d'argent en faisant des imports, c'est correct ?}

C'est exact. Elia est rémunérée pour la gestion du réseau indépendamment de l'origine géographique de l'électricité qui y transite. Que l'électricité soit produite en Belgique ou importée, cela ne modifie pas les revenus d'Elia. Il s'agit d'un mythe important à déconstruire : Elia n'a aucun intérêt économique à favoriser les importations.

\subsubsection{On est un couloir de passage pour d'autres pays, c'est correct ?}

Oui, la Belgique joue parfois un rôle de transit entre pays voisins~-- par exemple entre la France et l'Allemagne ou entre la France et les Pays-Bas. Ces flux de transit sont pris en compte dans les modèles d'adéquation. Ils constituent un avantage économique global : les interconnexions permettent d'exploiter la complémentarité des mix énergétiques européens (vent au nord, soleil au sud, nucléaire en France) et d'accéder à la production la moins chère à chaque instant, réduisant ainsi les coûts pour l'ensemble des consommateurs.

\subsubsection{Quelle est la complexité liée aux différents régulateurs ?}

Elia est soumise à quatre régulateurs~: la CREG (fédéral), la CWaPE (Wallonie), Brugel (Bruxelles) et la VNR anciennement VREG (Flandre). Cette multiplicité découle de la régionalisation belge~: le fédéral est compétent pour la sécurité d'approvisionnement, les grandes installations et l'offshore, tandis que les régions gèrent l'efficacité énergétique, l'aménagement du territoire, les permis et les infrastructures 30--70~kV. Elia doit donc élaborer quatre plans de développement~-- un fédéral et trois régionaux~-- selon des calendriers et des processus d'approbation différents. Cette architecture régulatoire complexifie notamment la mise en place de mécanismes de flexibilité et de raccordement cohérents à l'échelle nationale.

\subsubsection{Quels sont les différents timings de flexibilité qu'il faut avoir ? Quelles technologies agissent pour chacun de ces timings ?}

La flexibilité du système électrique s'organise selon plusieurs horizons temporels~:

\begin{itemize}
    \item \textbf{Secondes à minutes}~: les services auxiliaires (FCR, aFRR, mFRR) assurent le maintien de la fréquence à 50~Hz. Les batteries lithium-ion sont particulièrement adaptées à ces réponses rapides.
    \item \textbf{Quart d'heure à quelques heures}~: les batteries permettent la gestion des variations journalières de production renouvelable et de consommation.
    \item \textbf{Journée à semaine}~: la flexibilité de la demande entre en jeu~-- recharge des véhicules électriques à des horaires décalés, pilotage des pompes à chaleur via des buffers thermiques, modulation des consommations industrielles interruptibles. Environ 1~GW de gros consommateurs industriels réduisent déjà leur charge lorsque les prix de marché sont très élevés.
    \item \textbf{Plusieurs jours à saisonnier}~: les centrales thermiques (gaz, et à terme molécules défossilisées) constituent le dernier recours lors des épisodes de \textit{Dunkelflaute}. La flexibilité de la consommation ne peut pas couvrir ces périodes prolongées, car les industriels ne peuvent pas arrêter leurs équipements sur plusieurs jours.
\end{itemize}

Un problème émergent est celui de l'incompressibilité~: au printemps, lors des dimanches ensoleillés et venteux (notamment en mai), la production renouvelable peut excéder la consommation, imposant de réduire voire d'arrêter temporairement certaines installations photovoltaïques ou éoliennes. Elia publie chaque année un \textit{Summer Outlook} pour anticiper ces situations.

\subsubsection{Comment est-ce que les acteurs industriels reçoivent les demandes de consommation flexible ?}

Les connexions flexibles visent à raccorder plus rapidement certains clients en leur attribuant une puissance disponible la plupart du temps, mais non garantie à 100~\%~: pendant un certain pourcentage du temps, ils peuvent être invités à réduire leur soutirage. Les industriels dont les procédés sont peu flexibles trouvent parfois une solution dans le \textit{fuel switching}~: ils électrifient certains usages (e-boilers) tout en conservant leur chaudière au gaz, et basculent vers l'électricité lorsque son prix est très bas (voire négatif), avant de revenir au gaz sinon. Cette approche est économiquement intéressante pour l'industriel et constitue une aide précieuse pour le réseau. D'autres solutions incluent des buffers thermiques ou des stockages locaux. La CREG et les régulateurs sont impliqués dans la définition des règles encadrant ces mécanismes, dans le cadre plus large d'une réflexion sur les couloirs de nage par usage et sur l'évolution des processus de raccordement.